\chapter{Congruences}

\begin{center}\itshape
What are the last two digits of $3^{3^{3^3}}$, a number with more than $3\times10^{12}$ digits?\\
Arithmetic modulo $n$ answers in four lines.
\end{center}

\section{The ring \texorpdfstring{$\Z/n\Z$}{Z/nZ}}

\begin{reminder}[congruences]
For $n\ge1$, $a\equiv b\pmod n$ means $n\mid a-b$. Congruence is compatible with $+$ and $\times$, so the classes $\bar a=a+n\Z$ form a ring $\Z/n\Z$ with $n$ elements.
\end{reminder}

\begin{proposition}\label{prop:units}
$\bar a$ is invertible in $\Z/n\Z$ if and only if $\gcd(a,n)=1$. Hence $\Z/n\Z$ is a field if and only if $n$ is prime, and the group of units $(\Z/n\Z)^\times$ has $\varphi(n)$ elements.
\end{proposition}
\begin{proof}
$\bar a\bar u=\bar1$ means $au+nv=1$ for some $v$, which is B\'ezout's characterisation of $\gcd(a,n)=1$. The inverse is computed by the extended Euclidean algorithm.
\end{proof}

\begin{theorem}[Euler, Fermat]\label{thm:eulerthm}
If $\gcd(a,n)=1$ then $a^{\varphi(n)}\equiv1\pmod n$. In particular, for $p$ prime and $p\nmid a$, $a^{p-1}\equiv1\pmod p$.
\end{theorem}
\begin{proof}
Multiplication by $\bar a$ is a bijection of the finite group $(\Z/n\Z)^\times$ (its inverse is multiplication by $\bar a^{-1}$). So the product of all units $x$ equals the product of all $\bar ax$, which is $\bar a^{\varphi(n)}$ times the same product. Cancel that unit.
\end{proof}

\begin{theorem}[Wilson]\label{thm:wilson}
For $p$ prime, $(p-1)!\equiv-1\pmod p$.
\end{theorem}
\begin{proof}
In the field $\Z/p\Z$, $x^2=1$ means $(x-1)(x+1)=0$, so $x=\pm1$: only $1$ and $-1$ are their own inverses. The other non-zero classes pair off with their inverses, each pair having product $1$. So $(p-1)!\equiv1\cdot(-1)=-1$.
\end{proof}

\section{The Chinese remainder theorem}

\begin{theorem}[Chinese remainder theorem]\label{thm:crt}
If $m,n\ge1$ are coprime, the map $\Z/mn\Z\to\Z/m\Z\times\Z/n\Z$, $x\mapsto(x\bmod m,\ x\bmod n)$ is a ring isomorphism. Explicitly, if $mu+nv=1$, the solution of $x\equiv a\pmod m$, $x\equiv b\pmod n$ is $x\equiv anv+bmu\pmod{mn}$.
\end{theorem}
\begin{proof}
The map is a well-defined ring morphism. It is injective: if $m\mid x$ and $n\mid x$ with $m,n$ coprime then $mn\mid x$ (Gauss). Both sides have $mn$ elements, so it is bijective. For the formula: $nv\equiv1\pmod m$ and $nv\equiv0\pmod n$, and symmetrically for $mu$.
\end{proof}

Units correspond to pairs of units, so $\varphi(mn)=\varphi(m)\varphi(n)$ for coprime $m,n$. With $\varphi(p^k)=p^k-p^{k-1}$ (the non-multiples of $p$), this recovers $\varphi(n)=n\prod_{p\mid n}\left(1-\frac1p\right)$ from Tome~0.

\begin{proposition}[linear congruences]
$ax\equiv b\pmod n$ has a solution if and only if $d=\gcd(a,n)$ divides $b$, and then it has exactly $d$ solutions modulo $n$.
\end{proposition}
\begin{proof}
Solvability means $b\in a\Z+n\Z=d\Z$. If $d\mid b$, divide through: $\frac adx\equiv\frac bd\pmod{\frac nd}$ has a unique solution $x_0$ modulo $\frac nd$, since $\frac ad$ is a unit there; it gives $d$ classes modulo $n$.
\end{proof}

\section{Lifting modulo prime powers}

\begin{theorem}[Hensel's lemma]\label{thm:hensel}
Let $f\in\Z[X]$, $p$ prime, $k\ge1$, and $a\in\Z$ with $f(a)\equiv0\pmod{p^k}$ and $f'(a)\not\equiv0\pmod p$. Then there is a unique $t\in\{0,\dots,p-1\}$ such that $f(a+tp^k)\equiv0\pmod{p^{k+1}}$.
\end{theorem}
\begin{proof}
Taylor's formula for polynomials gives $f(a+h)=f(a)+hf'(a)+h^2g(a,h)$ with $g$ having integer coefficients. With $h=tp^k$ and $2k\ge k+1$:
$f(a+tp^k)\equiv f(a)+tp^kf'(a)\pmod{p^{k+1}}$. Writing $f(a)=cp^k$, the condition becomes $c+tf'(a)\equiv0\pmod p$, which has exactly one solution $t$ modulo $p$ since $f'(a)$ is invertible modulo $p$.
\end{proof}

Iterating, a simple root modulo $p$ lifts to a unique root modulo every $p^k$. This is Newton's method in the $p$-adic world (Chapter~12).

\section{Exercises}

\begin{exercise}[$\star$]\label{ex:last2}
Find the last two digits of $7^{2026}$.
\end{exercise}

\begin{exercise}[$\star$]\label{ex:sunzi}
\emph{Sunzi (3rd century).} Find the integers that leave remainders $2$, $3$, $2$ when divided by $3$, $5$, $7$.
\end{exercise}

\begin{exercise}[$\star$]\label{ex:digits}
Show that an integer is congruent modulo $9$ to the sum of its digits, and modulo $11$ to the alternating sum of its digits (starting from the units).
\end{exercise}

\begin{exercise}[$\star\star$]\label{ex:sqrt1}
\begin{enumerate}[label=(\alph*)]
\item Let $p$ be an odd prime and $k\ge1$. Show that the solutions of $x^2\equiv1\pmod{p^k}$ are $x\equiv\pm1$.
\item Let $n$ be odd with $r$ distinct prime factors. How many solutions does $x^2\equiv1\pmod n$ have?
\end{enumerate}
\end{exercise}

\begin{exercise}[$\star\star$]\label{ex:wilsonconv}
Show that if $n>4$ is composite then $(n-1)!\equiv0\pmod n$. Deduce that $n>1$ is prime if and only if $(n-1)!\equiv-1\pmod n$.
\end{exercise}

\begin{exercise}[$\star\star$]\label{ex:hensel}
Find $x\in\{0,\dots,124\}$ with $x^2\equiv-1\pmod{125}$, by lifting a solution modulo $5$.
\end{exercise}

\begin{exercise}[$\star\star$]\label{ex:tower}
Show that the last two digits of $3^{3^{3^3}}$ are $87$.
\end{exercise}

\begin{exercise}[$\star\star$]\label{ex:korselt}
\emph{Korselt's criterion, one direction.} Let $n$ be squarefree and composite, and suppose $p-1\mid n-1$ for every prime $p\mid n$. Show that $a^n\equiv a\pmod n$ for every integer $a$. Check that $561=3\cdot11\cdot17$ qualifies: it is a Carmichael number (Tome~0, Chapter~1 notebook).
\end{exercise}

\begin{exercise}[$\star\star\star$]\label{ex:lucas}
\emph{Lucas' theorem.} Let $p$ be prime, and write $n=\sum_in_ip^i$ and $k=\sum_ik_ip^i$ in base $p$. Show that
\[
\binom nk\equiv\prod_i\binom{n_i}{k_i}\pmod p,
\]
with $\binom{n_i}{k_i}=0$ when $k_i>n_i$. \emph{Hint: in $(\Z/p\Z)[X]$, $(1+X)^p=1+X^p$.} Deduce that $\binom nk$ is odd if and only if each binary digit of $k$ is at most the corresponding digit of $n$, which explains the Sierpi\'nski triangle of Tome~0.
\end{exercise}

\begin{beauty}
The Chinese remainder theorem says that arithmetic modulo $n$ is arithmetic modulo each prime power of $n$, independently. It makes RSA decryption four times faster (notebook), and in Tome~2 it becomes a statement about products of rings. Hensel's lemma is the door to the $p$-adic numbers of Chapter~12, where $\sqrt{-1}$ exists as a genuine $5$-adic number, the limit of $2,7,57,\dots$ from Exercise~\ref{ex:hensel}.
\end{beauty}

\section{Solutions}

\begin{solution}{ex:last2}
$7^2=49$ and $7^4=2401\equiv1\pmod{100}$. As $2026=4\cdot506+2$, $7^{2026}\equiv7^2=49\pmod{100}$.
\end{solution}

\begin{solution}{ex:sunzi}
Modulo $15$: $x\equiv2\pmod3$ and $x\equiv3\pmod5$ give $x\equiv8\pmod{15}$ (try $3,8$). Then $x=8+15t\equiv2\pmod7$ means $1+t\equiv2$, so $t\equiv1$ and $x\equiv23\pmod{105}$. The answers are $23+105k$.
\end{solution}

\begin{solution}{ex:digits}
If $n=\sum d_i10^i$ then $10\equiv1\pmod9$ gives $n\equiv\sum d_i$, and $10\equiv-1\pmod{11}$ gives $n\equiv\sum(-1)^id_i$.
\end{solution}

\begin{solution}{ex:sqrt1}
(a) $p^k\mid(x-1)(x+1)$, and $p$ cannot divide both factors (it would divide their difference $2$). So $p^k$ divides one of them.

(b) Write $n=p_1^{a_1}\cdots p_r^{a_r}$. By the Chinese remainder theorem (applied repeatedly), solving modulo $n$ amounts to solving modulo each $p_i^{a_i}$ independently, with $2$ solutions each by (a). So there are $2^r$ solutions.
\end{solution}

\begin{solution}{ex:wilsonconv}
Let $n=ab$ with $1<a\le b<n$. If $a<b$, both $a$ and $b$ appear among $1,\dots,n-1$, so $ab=n$ divides $(n-1)!$. If $a=b$, then $n=a^2$ with $a>2$ (as $n>4$), so $a<2a<a^2=n$ and $a\cdot2a=2n$ divides $(n-1)!$. For the equivalence: primes satisfy Wilson; $n=4$ gives $3!\equiv2$; composite $n>4$ gives $0\not\equiv-1$.
\end{solution}

\begin{solution}{ex:hensel}
Modulo $5$, $2^2=4\equiv-1$; with $f=X^2+1$, $f'(2)=4\not\equiv0$. Lift to $25$: $(2+5t)^2+1\equiv5+20t\pmod{25}$, so $1+4t\equiv0\pmod5$, $t=1$, $x=7$ (and $50\equiv0$). Lift to $125$: $(7+25t)^2+1\equiv50+350t\pmod{125}$, so $2+14t\equiv0\pmod5$, that is $4t\equiv3$, $t=2$, $x=57$. Check: $57^2+1=3250=26\cdot125$.
\end{solution}

\begin{solution}{ex:tower}
We need $3^E\bmod100$ with $E=3^{27}$. Since $\varphi(100)=40$ and $\gcd(3,100)=1$, $3^E$ depends only on $E\bmod40$. Now $3^4=81\equiv1\pmod{40}$, so $E=3^{27}=(3^4)^6\cdot3^3\equiv27\pmod{40}$. Hence $3^E\equiv3^{27}\pmod{100}$. Finally $3^{10}=59049\equiv49$, $3^{20}\equiv49^2=2401\equiv1$, so $3^{27}\equiv3^7=2187\equiv87\pmod{100}$.
\end{solution}

\begin{solution}{ex:korselt}
Let $p\mid n$ and write $n-1=(p-1)m$. If $p\nmid a$, Fermat gives $a^{n-1}=(a^{p-1})^m\equiv1$, so $a^n\equiv a\pmod p$; if $p\mid a$, both sides are $0$ modulo $p$. So every prime factor of $n$ divides $a^n-a$, and since $n$ is squarefree, their product $n$ does (Chinese remainder theorem). For $561$: $560$ is divisible by $2$, $10$ and $16$.
\end{solution}

\begin{solution}{ex:lucas}
In $(\Z/p\Z)[X]$, $(1+X)^p=1+X^p$ because $p\mid\binom pj$ for $0<j<p$ (Tome~0, Exercise~1.6); by induction $(1+X)^{p^i}=1+X^{p^i}$. Hence
\[
(1+X)^n=\prod_i\left(1+X^{p^i}\right)^{n_i}=\prod_i\sum_{j=0}^{n_i}\binom{n_i}jX^{jp^i}.
\]
Expanding, each exponent $k$ arises from exactly one choice of digits $j=k_i\le n_i$ (uniqueness of base-$p$ expansion, since all $j\le n_i<p$), with coefficient $\prod_i\binom{n_i}{k_i}$; if some $k_i>n_i$, $X^k$ does not appear. Compare with the coefficient $\binom nk$ of $X^k$. For $p=2$, $\binom{n_i}{k_i}$ is $0$ exactly when $k_i=1$, $n_i=0$.
\end{solution}
